\documentclass[11pt]{article}
\usepackage[a4paper,margin=25mm]{geometry}
\usepackage{amsmath,amssymb,amsthm,mathtools,booktabs,hyperref,xurl}
\setlength{\emergencystretch}{2em}
\newtheorem{theorem}{Theorem}
\newtheorem{lemma}[theorem]{Lemma}
\newtheorem{proposition}[theorem]{Proposition}
\newtheorem{corollary}[theorem]{Corollary}
\newcommand{\Z}{\mathbb Z}\newcommand{\Q}{\mathbb Q}
\title{A short Pell replacement for the fixed Jones power}
\author{\(\mathbb Z\mathbb Q\) Diophantine Project}
\date{23 November 2026}
\begin{document}
\maketitle

\begin{abstract}
We replace the long fixed-power circuit for \(q=b^{5^{60}}\) by a
constant-size quadratic Pell certificate with two additional integer
inputs. A specially chosen auxiliary parameter first makes the
geometric quotient \((q^4-1)/(b^5-1)\) integral without knowing the
exponent. The retained binomial block can then be reconstructed in a
noncircular order. Its explicit growth estimates recover the intended
exponent and exclude the unwanted signed Pell branch. Only two new
positivity requirements are needed. Complete coefficient tables using
thirteen integer inputs have 704 variables and 595 quadratic residuals
for the ordinary-\(abc\) interface, or 979 variables and 829 residuals
for the general interface; their sums of squares have degree four.
The local source equivalence
requires no additional arithmetic or coding hypothesis. Its application
to rational undecidability retains the stated hypotheses of the
chosen integer interface and the previously cited Jones coding theorem.
\end{abstract}

\section{The exact replacement}

Set \(L=5^{60}\) and \(T_b=b^5-1\). Retain the corrected source,
its one-quotient third Pell certificate, and its half-trace branch
square from \cite{CRT,Halftrace}, but delete the fixed-power circuit
for \(q=b^L\). The existing integer inputs and all other equations
are retained. Restrict the integer instance to the cofinite range
\[
                              x>32z.
\]
Introduce two integer inputs \(\delta,\alpha_*\), and define
\begin{align}
 A_0&=(b^9+8b^5+b-8)/2,& A&=A_0+5T_bs,\label{A}\\
 m_*&=2bA-b^2-1,& C_*&=L+\delta(A-1).\label{mC}
\end{align}
Require a rational \(D_*\) to satisfy
\begin{align}
 D_*^2-(A^2-1)C_*^2&=1,\label{norm}\\
 q&=D_*-(A-b)C_*-\alpha_*m_*,\label{q}
\end{align}
and add precisely the two strict inequalities
\begin{equation}\label{ineq}
                         q>b^5,\qquad C_*^2<c^2.
\end{equation}
Here \(s,c\) are the retained integer inputs of the binomial block.
The new \(\alpha_*\) is unrelated to the retained positive source
slack \(\alpha\). The old guard \(b^5-2z>0\), with integral
\(b\) and \(z\ge1\), still gives \(b\ge2\).

\section{Integral reconstruction before exponent recovery}

\begin{lemma}\label{parameter}
For integers \(b\ge2\) and \(s\), the quantities in
\eqref{A}--\eqref{mC} satisfy
\[
 A_0,A,m_*\in\Z,\quad A\equiv1\pmod5,\quad
 m_*=T_b(b^5+8b+1+10bs),\quad |A|\ge2,\quad 0<A_0<b^{16}.
\]
\end{lemma}
\begin{proof}
The powers \(b^9\) and \(b\) have the same parity.
Modulo five, \(b^9\equiv b^5\equiv b\), so
\(2A_0\equiv10b-8\equiv2\), whence \(A_0\equiv A\equiv1\).
Direct expansion proves the displayed factorization of \(m_*\).
If \(A=1\), it would imply \(T_b\mid-(b-1)^2\), impossible
because \(0<(b-1)^2<T_b\). The congruence modulo five excludes
\(A=0,-1\), proving \(|A|\ge2\).
Finally \(A_0>0\), and \(8b^5\le b^8\), \(b<b^8\) give
\[
       2A_0<b^9+2b^8\le2b^9<2b^{16}.
\]
\end{proof}

Use the integer polynomial Pell sequences \(X_j(A),Y_j(A)\),
defined by \(X_0=1,Y_0=0,X_1=A,Y_1=1\) and the recurrence
\(Z_{j+2}=2AZ_{j+1}-Z_j\). For negative parameters,
\[
 X_j(-A)=(-1)^jX_j(A),\qquad
 Y_j(-A)=(-1)^{j-1}Y_j(A).
\]
Integral Pell classification at \(|A|\ge2\) is proved in
\cite{Index}.

\begin{lemma}[The geometric quotient is already integral]\label{lambda}
With integer \(b,s,\delta,\alpha_*\), \(b\ge2\), and
\eqref{A}--\eqref{q}, the quantities \(C_*,D_*,q\) are integers
and \(q\equiv\pm1\pmod{T_b}\). Hence the retained equation
\(\lambda T_b=q^4-1\) gives \(\lambda\in\Z\). Under
\(q>b^5\), this integer is positive.
\end{lemma}
\begin{proof}
First \(C_*\in\Z\). The norm makes \(D_*^2\) integral, so
rationality makes \(D_*\) integral; then \eqref{q} makes \(q\)
integral. Signed Pell classification gives some \(\nu\ge0\) and
independent signs with
\[
                         C_*=\pm Y_\nu(A),\qquad D_*=\pm X_\nu(A).
\]
The zero index is allowed here. As \(A\equiv1\pmod5\), the
recurrence gives \(Y_\nu(A)\equiv\nu\pmod5\). Equation
\eqref{mC} makes \(5\mid C_*\), and thus \(5\mid\nu\).

Put \(R_j=X_j(A)-(A-b)Y_j(A)\). The addition and norm identities
give
\begin{equation}\label{rec}
\begin{split}
 R_0&=1,\qquad R_{j+1}=bR_j+m_*Y_j(A),\\
 R_j\bigl(X_j(A)+(A-b)Y_j(A)\bigr)&=1+m_*Y_j(A)^2.
\end{split}
\end{equation}
Modulo \(T_b\), these give \(R_j\equiv b^j\) and the other
factor congruent to \(b^{-j}\), since \(\gcd(b,T_b)=1\).
The four signs of \(C_*,D_*\) therefore make
\(D_*-(A-b)C_*\) congruent to one of
\(\pm b^\nu,\pm b^{-\nu}\). All are \(\pm1\) because
\(5\mid\nu\) and \(b^5\equiv1\pmod{T_b}\).
Now use \(T_b\mid m_*\) and \eqref{q}. The conclusions about
\(\lambda\) follow because \(T_b>0\) and \(q>b^5\).
\end{proof}

This recovers both \(q\) and \(\lambda\) integrally before proving
\(q=b^L\), before knowing \(s>0\), and before invoking the source's
binomial Pell indices.

\section{The cofinite tail restores prefix positivity}

Let \(n=q^{16}\), \(M=n^2=q^{32}\), and retain the source formula
\[
                         r=S(M-n)+T(M-1),
\]
with \(S,T\) as in the corrected Jones coding prefix.

\begin{lemma}\label{prefix}
The retained source guards, \(x>32z\), positive integral
\(\lambda\), and integral \(q>b^5\) imply
\[
                        S,T\in\Z_{>0},\qquad
                        r\ge M>b^{160}\ge2^{160}>L.
\]
\end{lemma}
\begin{proof}
The positive integer inputs \(g,t,m\) give
\(e=y+m(b^5-2z)>0\) and
\(\ell=u+t(b^5-2z)>0\). The retained positive slack says
\[
                       e\ell g^2<(b-xy)q^2.
\]
Thus \(b>xy\ge x>32z\), and \(x\le b-1\).
Since \(2z<b\), one has \(e,\ell\ge b^5-b+2\ge b^4\),
and hence \(g<q/b^3\). The quantity \(c_0=1+xb^5+g\)
satisfies
\[
                  0<c_0<b^6+q/b^3<(b+b^{-3})q<2bq,
\]
using \(q>b^5\). Consequently
\[
 b^5(1+q^4)-2zc_0^4
       >b^5+b^4q^4(b-32z)>0.
\]
The original prefix expression
\[
 S=g+eq^3+\ell q^5+
       [2ec_0^4+\lambda\{b^5(1+q^4)-2zc_0^4\}]q^4
\]
is positive. Also \(\ell<bq^2\), so
\[
 T=q^3-(b-1)\ell+(b^5-2z)\lambda q^3+(b^5-2)q^5>0,
\]
because \((b^5-2)q^3>b(b-1)\). All quantities are integral.
Therefore \(S,T\ge1\) and
\(r\ge2M-n-1\ge M\) for \(n\ge2\).
Finally \(M=q^{32}>b^{160}\), while
\(L=5^{60}<2^{140}<2^{160}\).
\end{proof}

Every prerequisite of the retained signed binomial argument now has
its established noncircular reconstruction. In particular \(r,a,p,k\)
are integral; the half-trace square excludes \(s=0\); the strict
small-Pell inequality gives \(c\ne0\); the reciprocal row then
legitimizes the integer \(i\ne0\). Also \(M>2^{160}\) and
\(2\le b<q\le2M\). The local theorems of
\cite{Signed,CRT,Halftrace} therefore give
\begin{equation}\label{canonical}
\begin{split}
 &w,s,c>0,\qquad bw=2^{2r+1},\\
 &c=Y_{2r+1}(a),\qquad d=X_{2r+1}(a),\qquad
 M\mid\binom{2r}{r}.
\end{split}
\end{equation}
These local theorems require the stated integral parameters and
bounds, not the equality \(q=b^L\) now being recovered.

\section{Growth identifies the auxiliary index and its sign}

Put \(U=wM\), \(Z=sM\), \(B=2r+1\), so \(a=rsM(U+1)\).
The signed binomial proof supplies \(Z>\tfrac78(U-2)^r\).

\begin{lemma}\label{growth}
Every proposed solution satisfies
\[
 U\ge4^r,\quad s>U^{r-2},\quad a<s^2,\quad
 s>L+6r+2,\quad A-1>L+6r,\quad m_*>s>0.
\]
\end{lemma}
\begin{proof}
From \(bw=2^{2r+1}\) and \(b\le2M\) one gets \(U\ge4^r\),
so \(2r/U\le1/2\). Bernoulli's inequality gives
\((1-2/U)^r\ge1-2r/U\ge1/2\). As \(M\le U\),
\[
 s>\frac7{16}\frac{U^r}{M}\ge\frac7{16}U^{r-1}>U^{r-2}.
\]
Also \(2r<U\) and \(r>2^{160}\), so
\[
                       a<2rsU^2<sU^3<s^2.
\]
Since \(r>L\), the elementary inequality \(4^r>7r+2\) for
\(r\ge3\) gives \(s>U\ge4^r>L+6r+2\).
Now \(A=A_0+5T_bs>s\) because \(s>0,A_0>0,T_b\ge31\).
Finally
\[
                   m_*>10bT_bs-b^2-1>s;
\]
the last step follows from \(s\ge1,T_b\ge b\), and
\(10b^2-1>b^2+1\).
\end{proof}

\begin{lemma}\label{index}
The index in Lemma~\ref{lambda} is \(\nu=L\), and
\(C_*=Y_L(A)>0\).
\end{lemma}
\begin{proof}
After Lemma~\ref{growth}, \(A>0\).
If \(C_*=0\), then \(A-1\mid L\) by \eqref{mC},
contradicting \(A-1>L\). Thus \(\nu\ge1\).
Elementary Pell growth and \eqref{ineq} give
\[
 A^{\nu-1}\le Y_\nu(A)=|C_*|
       <c=Y_{2r+1}(a)<(2a)^{2r}<s^{6r}.
\]
Since \(A>s\), this implies \(\nu\le6r\).
Reduction of \(C_*=L+\delta(A-1)\) modulo \(A-1\) gives
\(\nu\equiv L\) or \(-L\), according to its sign.
The bounds \(0<\nu\le6r\) and \(A-1>L+6r\) allow only
\(\nu=L\) with the positive sign. In the negative case,
\(A-1\) would divide the positive integer \(\nu+L<A-1\).
\end{proof}

\begin{theorem}[Exact power recovery and completeness]\label{main}
Every rational solution of the replacement satisfies \(q=b^L\).
Conversely every old corrected source solution with \(x>32z\)
extends to a replacement solution with \(\delta,\alpha_*>0\).
The two sources have the same coding-prefix and instance projection.
\end{theorem}
\begin{proof}
By Lemma~\ref{index}, \(C_*=Y_L(A)>0\) and
\(D_*=\pm X_L(A)\). Set \(K=b^L\).
The identities \eqref{rec}, now modulo \(m_*>0\), give
\[
 \begin{array}{ll}
 D_*=X_L(A):&q\equiv K\pmod{m_*},\\
 D_*=-X_L(A):&qK\equiv-1\pmod{m_*}.
 \end{array}
\]
The inverse congruence is legitimate since
\(m_*\equiv-1\pmod b\). As \(b<q=M^{1/32}\le U^{1/32}\),
\[
                  qK\le U^{(L+1)/32}\le U^{(6r+1)/32}.
\]
For \(r\ge4\), one has \((6r+1)/32+1\le r-2\). Since \(U>2\),
\[
 0<qK+1\le2qK\le2U^{(6r+1)/32}<U^{r-2}<s<m_*.
\]
Thus the negative-trace case is impossible. In the positive case,
\(|q-K|<qK+1<m_*\), so the congruence forces \(q=K=b^L\).

Conversely take an old corrected source solution.
It has \(q=b^L>b^5\), the canonical values \eqref{canonical},
and all the same growth bounds. Since
\(s>U^{r-2}>b^{16}>A_0\),
\[
                       A<(1+5T_b)s<5b^5s<a,
\]
where \(a\ge Ms>5b^5s\). Also \(L<r<2r+1\).
Set \(C_*=Y_L(A)\), \(D_*=X_L(A)\). Monotonicity in the
positive parameter and index gives
\(0<C_*<Y_{2r+1}(a)=c\), proving the second bound in
\eqref{ineq}. The recurrence modulo \(A-1\) makes
\[
                      \delta=\frac{Y_L(A)-L}{A-1}
\]
an integer, positive because \(Y_L(A)>L\).
Iteration of \eqref{rec} gives
\[
 R_L-b^L=m_*\sum_{j=0}^{L-1}b^{L-1-j}Y_j(A).
\]
The sum is a positive integer because \(L\ge2\).
Choose it as \(\alpha_*\); this proves \eqref{q}.
Every other coordinate is retained. Conversely the already proved
power equality restores the deleted circuit by its defining values.
\end{proof}

Deleting the finitely many old instance values \(x\le32z\) preserves
undecidability: a decision procedure for a cofinite tail, together
with a fixed finite lookup table, would decide the original set.

\section{Quadratic compilation}

Apart from its two integer inputs, the replacement uses six new
rational registers
\[
                         b_4,A,C_*,D_*,m_*,A_C.
\]
For the ordinary-\(abc\) interface, the existing input square is
\(b_2=b^2\), and the new equations are
\begin{align*}
 b_4&=b_2^2,& b^5&=b\,b_4,\\
 2A&=b_4b^5+8b^5+b-8+10(b^5-1)s,&&\\
 C_*&=L+\delta(A-1),& A_C&=AC_*,\\
 D_*^2-A_C^2+C_*^2&=1,&m_*&=2bA-b^2-1,\\
 q&=D_*-A_C+bC_*-\alpha_*m_* .&&
\end{align*}
Here \(b^5\) is a retained coordinate, so all residuals are quadratic.
The general interface stores \(b^2/\rho\) and \(1/\rho^2\);
there the first row is instead
\((1/\rho^2)b_4=(b^2/\rho)^2\).
The retained global gauge guard already gives \(\rho\ne0\), so
this has the same unique value \(b_4=b^4\).
No new reciprocal is introduced. The row for \(b^5\) replaces
the one using the discarded \(b^3\) register.

Each strict integer inequality in \eqref{ineq} uses three rational
triangular-number coordinates and one quadratic residual. The private
positive representatives required by these rows exist for every
positive integer value. The new integer inputs \(\delta,\alpha_*\)
are positive in the completeness proof, so the nonzero-input scope
of the ordinary-\(abc\) simultaneous certificate is respected.

The source proof above is independent of the choice of integer
interface. A complete quartic record must additionally instantiate
its thirteen integer certificates and shared contacts. The exact
coefficient builders and independent audits accompanying this paper
perform that step. The resulting arithmetic transfer retains all
premises of the chosen interface and the same cited Jones coding
theorem; no new source universality assumption is introduced.

\subsection*{The complete thirteen-input ordinary-\(abc\) family}

The existing fixed-offset theorem in
\nolinkurl{undecidability/papers/abc_fixed_contact_offsets_v93.tex}
is stated for an arbitrary fixed number \(k\) of nonzero integer
inputs. After common-reference and common-square sharing, its count,
including the inputs, is \(48k+8\) coordinates and \(41k+5\)
quadratic residuals. Its soundness applies to every rational zero;
its simultaneous completeness supplies every tuple in
\((\Z\setminus\{0\})^k\), under the same retained arithmetic
premises and the same sufficiently large fixed odd coefficient.

The eleven retained inputs use the old dyadic offsets
\[
       0,1,2,3,4,5,6,7,8,10,11\quad\hbox{divided by }16.
\]
Assign \(\delta\) offset \(9/16\) and \(\alpha_*\) offset
\(12/16\). These are precisely thirteen distinct cosets modulo
\(\Z_2\). For each nonzero integer input, the individual contact
positions are in \(\Z_2\); subtracting these offsets therefore
separates all the contact lists. The finite-family supplier applies
to this larger fixed union, with the previous common reference,
the same six shared reference coordinates, and the same common
contact square. Its temporary finite avoidance is enlarged by the
new positions, their differences, and denominators exactly as in
the fixed-offset theorem. No uniform bound on an input-dependent
annihilation order is needed, and no moving reference-dependent
offset is introduced.

For soundness, each of the thirteen blocks recovers a complete copy
of the integer-defining interface. Thus all old inputs and the new
\(\delta,\alpha_*\) are integers, which is exactly the integrality
premise used in Lemma~\ref{lambda}. For completeness, start with a
positive corrected source witness. Theorem~\ref{main} supplies
\(\delta,\alpha_*>0\), while the retained positive reconstruction
supplies nonzero values of the other eleven inputs. The
thirteen-input fixed-offset theorem therefore supplies all the
rational integer certificates simultaneously. This is an application
of the existing arbitrary-fixed-\(k\) theorem, with no extension of
its arithmetic hypotheses.

The coefficient construction copies the old \(o\) block, whose
input had offset \(9/16\), twice. That block has 47 private
auxiliary coordinates and 41 residuals after sharing. One copy is
renamed for \(\delta\) without changing its offset. To obtain the
\(\alpha_*\) copy, write \(z\) and \(q_z=z^2\) for the shared
contact coordinates and make the simultaneous affine substitution
\[
             z\longmapsto z+\frac3{16},\qquad
             q_z\longmapsto q_z+\frac38z+\frac9{256}.
\]
The shared square equation makes the second substituted value exactly
\((z+3/16)^2\). Hence the physical offset changes from \(9/16\)
to \(12/16\), in every affected residual, including squared
occurrences of a contact invariant. All substitutions are affine
in the stored coordinates, so each residual remains quadratic.
Constant denominators are cleared and only nonzero rational row
scalings are used. All six common reference coordinates are kept
unchanged. A new rational zero consequently recovers the old
single-block equations; conversely the chosen thirteen physical
contact witnesses satisfy their copied equations.

The two new integer inputs and their private certificates add
\(2+2\cdot47=96\) coordinates and \(82\) residuals to the source
core. The complete table has 704 coordinates and 595 quadratic
residuals. Its sole polynomial is the sum of their squares, with
integer coefficients in the fixed parameters and the instance.
It has degree exactly four because an unchanged input-square row
has a nonzero quadratic homogeneous part. Over \(\Q\), its zero
set is exactly the simultaneous zero set of all 595 residuals.
The file
\nolinkurl{undecidability/pell_power_v97/pell_power_abc_quadrics.json}
specifies every coordinate and coefficient; the independent checker
reconstructs all 82 copied rows, the source-core rows, every
unchanged predecessor row, the guards, and the thirteen dyadic
cosets. The source-coding theorem and all retained ordinary-\(abc\)
interface premises remain explicit qualifications.

\section{Thirteen simultaneous general integer certificates}
\label{sec:thirteen-general-certificates}

The power replacement adds the integer inputs
\(\alpha_{\rm pow}\) and \(\delta_{\rm pow}\).
They are the literal coordinates \texttt{power\_alpha} and
\texttt{power\_delta}. The preceding general quartic already has
eleven integer certificates, with one common reference, common
numerator and scalar point data, a global contact gauge, and the
shared affine progression of first contact arguments. Its ordinary
integer predicates accept every integer, including zero, so no extra
positivity or nonzero requirement is imposed on the two new inputs.

\begin{theorem}[The thirteen-block general transfer]
Under the retained arithmetic premises of the general integer
interface, there is one sufficiently large fixed general defining
parameter \(C_{\rm parameter}=149^N\) for which the complete power
replacement has the required undecidable instance projection on the
cofinite coding range \(x>32z\). The complete table has 979 rational
variables and 829 quadratic residuals. Their sum of squares has degree
exactly four. No effective value of the sufficient \(N\) is asserted.
\end{theorem}
\begin{proof}
First instantiate the existing general simultaneous certificate with
the thirteen inputs
\[
 b,g,t,m,w,s,h,c,\gamma,j,i,\delta_{\rm pow},\alpha_{\rm pow}.
\]
The reciprocal chart of \(i\) is unchanged. The other twelve
certificates are ordinary integer-input charts. Retain the first
eleven contact-progression indices
\[
                    0,1,2,3,4,5,6,7,8,10,11
\]
and assign indices 12 and 13 to the new inputs, respectively.
These are distinct integer indices; no argument requires consecutive
indices.

There are two contact weights per input. Replace the previous sum of
22 weights by the sum of the 26 retained and new weights. At a
consumed prime belonging to any one of the thirteen blocks, the
global-gauge soundness proof uses that block's own tests, complementary
weight sum, and pole obstruction. Its valuation and parity arguments
do not depend on the presence of precisely eleven blocks. The same
finite-family height estimates therefore give integrality of all
thirteen distinguished inputs. The numerical constants now depend
on this fixed larger family. Choose \(N\) sufficiently large for
these constants and for the original scalar completeness threshold.
This is a single choice before the coding instance is varied. The
old eleven-block sufficient threshold is not claimed to remain
unchanged.

For completeness, fix a source witness with its thirteen integer
inputs. Apply the existing ordered joint scalar preparation and the
common numerator theorem to this finite tuple. Its difference points,
normalizers, and aggregate contact shifts are now fixed. Choose one
positive integer spacing larger than the diameter of their union of
contact lists. Translating list \(\ell\) by its distinct integer
index times that spacing separates the lists, by the same interval
argument as in the original progression theorem. The spacing depends
only on these fixed lists. In particular it does not depend on the
reference index subsequently adjusted at the comparison primes.

Apply the same retained finite-family prime supplier to the resulting
fixed rational positions, enlarging finite avoidance by all thirteen
inputs' fixed data. Its all-sign order conditions give the same common
reference and common-numerator congruences. At a selected prime,
only its own weight has negative odd valuation; every other weight
is integral. The global sum and its complementary sums consequently
have the required valuations. At foreign contacts the unchanged
rational-square parity argument excludes negative odd denominators.

At the prepared finite places and at infinity, with all translations
now fixed, every weight has its old asymptotic
\(q_\nu=\kappa\lambda^{-2}+O(\lambda^{-3})\). Thus
\[
 \rho=26\kappa\lambda^{-2}+O(\lambda^{-3}),\qquad
 \frac{(\rho-q_\nu)^2}{q_\nu}
       =25^2\kappa\lambda^{-2}+O(\lambda^{-3})\longrightarrow0.
\]
The argument divides by neither 26 nor 25. The source physical tests
remain bounded on their fixed neighborhoods; detector errors tend to
zero, and the actual native-pole arguments retain their prepared
local square classes and the old positive real limit. The supplier's
finite divisibility and real-slope clauses make these conditions hold
simultaneously. Apply the same native realization independently to
each actual pole input. The affine progression reconstruction then
fills the six common coefficient equations and every remaining
private chart row.

This is the existing finite-family arithmetic proof with thirteen
members. It introduces no new prime-supplier, elliptic, height, or
native-realization premise. Its source projection is the proved Pell
power replacement. Removing the fixed finite initial range
\(x\le32z\) does not change undecidability of the selected coding
set. The coefficient census and degree assertion are verified below.
\end{proof}

The retained family proofs are
\emph{One contact gauge in a combined quadratic system},
\nolinkurl{undecidability/papers/global_gauge_quartic_checkpoint.tex};
\emph{Common point coordinates in the quartic transfer},
\nolinkurl{undecidability/papers/common_points_quartic_checkpoint.tex};
and \emph{A common arithmetic progression of contact arguments},
\nolinkurl{undecidability/general_contact_v94/affine_contact_progressions.tex}.
All changes here concern the fixed cardinality of their finite family
and its complete coefficient implementation.

\paragraph{Exact implementation and census.}
The source-only replacement removes 137 power coordinates and adds
14 coordinates, including the two new integer inputs. It removes
138 source residuals and adds nine, while replacing the definition
of \(B_5\). Applied to the preceding 982-coordinate, 854-row
general system, this gives 859 coordinates and 725 rows before the
two new integer certificates are supplied.

Each new ordinary certificate contributes 60 auxiliary coordinates
and 52 rows. Its input is already counted in the new source, so it
is not counted again. The complete census is therefore
\[
               859+2\cdot60=979,\qquad
               725+2\cdot52=829.
\]
The new source core itself has 84 coordinates and 56 residuals.

The compiler clones the old ordinary \(o\)-input chart, whose first
contact used progression index 9. Its remaining private coordinates
and rows are copied with their input and prefix renamed. Graded
rescaling of the seven common progression coefficients changes index
9 to index 12 or 13. Constant denominators are cleared and affected
rows are made primitive. This is exactly the affine substitution
\[
 \begin{aligned}
 t_\ell&=t+\ell d,\\
 u_\ell&=u+2\ell v+\ell^2w,\\
 f_\ell&=f+\ell a_1+\ell^2a_2+\ell^3a_3+\ell^4a_4.
 \end{aligned}
\]
It adds no new common progression register or equation.

The two new weights from each block are added to the three rows
where the global sum was eliminated in the final reciprocal-\(i\)
contact. These are its pole-input definition, inverse row, and gate
row. Twelve new quadratic terms implement the four new weight
summands in those three rows. No other old integer-certificate row
is altered, and no existing gauge or nonzero guard is removed.

Every residual has coefficients in the integer polynomial ring of
the fixed coefficient parameters and has degree at most two in the
counted rational variables. At least one retained residual has a
nonzero quadratic homogeneous part. The sum of the squares of those
parts cannot vanish identically over the reals, so the final degree
is exactly four.

\paragraph{Independent verification.}
The complete coefficient table and certificate are
\begin{center}\small
\nolinkurl{undecidability/pell_power_v97/power_general_m163_quadrics.json}\\
\nolinkurl{undecidability/pell_power_v97/power_general_m163_certificate.json}.
\end{center}
The independent checker imports no builder. It specifies the nine
new source residuals independently, checks every removed-coordinate
occurrence, compares every unchanged row, and audits all twelve new
global-sum terms. To check the new private certificates it starts
from the earlier \emph{unprogressed} corrected \(o\)-input chart
and directly substitutes indices 12 and 13. This is independent of
the compiler's graded scaling of its already progressed template.
It checks all 104 private rows, retained guards, parameter domains,
support, degree, counts, and predecessor and output hashes. Run
\begin{verbatim}
python undecidability/pell_power_v97/build_power_general.py
python undecidability/pell_power_v97/verify_power_general.py
\end{verbatim}
The finite checks certify the complete algebraic implementation;
the quantified arithmetic completeness and source universality are
the proofs above and the retained external coding theorem.


\section{Exact local-lemma verification}

The independent script
\begin{verbatim}
python undecidability/pell_power_v97/verify_power_lemmas.py
\end{verbatim}
checks the parameter factorization, parity and congruence statements,
Pell recurrence and product identities, all signed congruence
branches, and the deliberately permitted preliminary index zero.
Its finite fixtures include negative auxiliary parameters. The
quantified arguments are the proofs above; finite fixtures do not
replace them. The complete-table checkers in the preceding two
sections additionally audit every coefficient change, shared-family
lift, count, guard, and hash.

\begin{thebibliography}{9}
\small\setlength{\itemsep}{2pt}
\bibitem{Index} \(\mathbb Z\mathbb Q\) Diophantine Project,
\emph{An exact three-Pell index certificate and nine fewer quadratic
coordinates}.\par
\nolinkurl{undecidability/papers/pell_index_optimization.tex}.
\bibitem{Signed} \(\mathbb Z\mathbb Q\) Diophantine Project,
\emph{Signed binomial witnesses force their own positivity and save
five quadratic coordinates}.\par
\nolinkurl{undecidability/source_repair_v93/signed_binomial_positivity.tex}.
\bibitem{CRT} \(\mathbb Z\mathbb Q\) Diophantine Project,
\emph{A single integer quotient for the third Pell certificate}.\par
\nolinkurl{undecidability/pell_crt_v95/one_quotient_pell.tex}.
\bibitem{Halftrace} \(\mathbb Z\mathbb Q\) Diophantine Project,
\emph{A Pell half-trace square absorbs a nonzero guard}.\par
\nolinkurl{undecidability/pell_halftrace_v96/halftrace_guard.tex}.
\end{thebibliography}
\end{document}
